% % % % % % % % % % % % % % % % % % % % % % % % % % % % % % % % % %
%  Devolutiva — Fase 1: Imersão
%  ETE Cícero Dias · Módulo 1 · Turma A · 2026.1
%  Profa. Samara Silva  ·  Recife, 1 de junho de 2026
% % % % % % % % % % % % % % % % % % % % % % % % % % % % % % % % % %

\documentclass[
  sigconf,
  language=portuguese
]{acmart}

% ── Desliga rodapé ACM ──────────────────────────────────────
\settopmatter{printacmref=false}
\setcopyright{none}
\acmYear{2026}

% ── Pacotes ─────────────────────────────────────────────────
\usepackage[portuguese]{babel}
\usepackage{microtype}
\usepackage{booktabs}
\usepackage{tabularx}
\usepackage{array}
\usepackage{colortbl}
\usepackage[dvipsnames,table]{xcolor}
\usepackage[inline]{enumitem}
\usepackage{graphicx}
\usepackage{textcomp}

% ── Cores ───────────────────────────────────────────────────
\definecolor{cgood}{HTML}{1E6B3A}
\definecolor{cwarn}{HTML}{8B4500}
\definecolor{cred}{HTML}{C0392B}
\definecolor{cmuted}{HTML}{666666}
\definecolor{crow}{HTML}{F2EFEB}

% ── Colunas customizadas ─────────────────────────────────────
\newcolumntype{L}[1]{>{\raggedright\arraybackslash}p{#1}}
\newcolumntype{C}[1]{>{\centering\arraybackslash}p{#1}}

% ── Comandos ──────────────────────────────────────────────────
\newcommand{\notadestaque}[2]{%
  \noindent\textbf{Nota:}\enspace%
  {\large\bfseries\color{cred}#1\,/\,#2\,pts}%
}
\newcommand{\statusok}[1]{{\small\bfseries\color{cgood}#1}}
\newcommand{\statuswarn}[1]{{\small\bfseries\color{cwarn}#1}}

% ── PT-BR ────────────────────────────────────────────────────
\AtBeginDocument{%
  \renewcommand{\abstractname}{Resumo}
  \renewcommand{\figurename}{Figura}
  \renewcommand{\tablename}{Tabela}
}

% ────────────────────────────────────────────────────────────
\begin{document}

% ── CAPA ────────────────────────────────────────────────────
\thispagestyle{empty}
\IfFileExists{capa.png}{%
  \includegraphics[width=\paperwidth,height=\paperheight]{capa.png}%
}{%
  \begin{center}\small\color{gray}[capa.png n\~ao encontrada]\end{center}%
}
\newpage

% ── METADADOS ────────────────────────────────────────────────
\title{Devolutiva de Avalia\c{c}\~ao --- Grupo 03 - Doação de Alimentos}

\author{Profa.\ Samara Silva}
\authornote{Grupo: \textbf{Grupo 03 - Doação de Alimentos} \textperiodcentered\ Per\'{i}odo: 22 de maio de 2026 \textperiodcentered\ Avaliado em: 1 de junho de 2026.}
\email{samarasilvia.educa@gmail.com}
\affiliation{%
  \institution{ETE C\'{i}cero Dias}
  \department{M\'{o}dulo 1 --- Turma A · 2026.1}
  \city{Recife}
  \state{Pernambuco}
  \country{Brasil}
}

\begin{abstract}
Devolutiva de avaliação do grupo \textbf{Grupo 03 - Doação de Alimentos}, referente à Fase~1 --- Imersão, no âmbito das disciplinas do Módulo~1 da ETE Cícero Dias. A nota total obtida foi de \textbf{2,47} pontos.
\end{abstract}

\maketitle

\section{Design Thinking}

\notadestaque{2,47}{10,00}

\subsection{Resumo dos Critérios}

\rowcolors{2}{crow}{white}
\begin{tabular}{@{} L{2.8cm} C{2.2cm} C{1.1cm} C{1.1cm} @{}}
\toprule
\textbf{Critério} & \textbf{Status} & \textbf{Nota} & \textbf{Máx.} \\
\midrule
Pesquisa Desk & \statuswarn{Completo c/ ressalvas} & \textbf{0,42} & 0,50 \\
Matriz de Alinhamento & \statuswarn{Faltou pouco c/ erros} & \textbf{0,30} & 0,50 \\
Matriz CSD & \statuswarn{Completo, inadequado} & \textbf{0,35} & 0,50 \\
Pesquisa Primária & \statuswarn{Completo, inadequado} & \textbf{0,70} & 1,00 \\
Relatório de Imersão & \statuswarn{Completo, inadequado} & \textbf{0,70} & 1,00 \\
\midrule
\textbf{Total} & & \textbf{\color{cred}2,47} & \textbf{3,50} \\
\bottomrule
\end{tabular}

\section{Comentários por Critério}

\subsection{Pesquisa Desk}

\noindent\statuswarn{0,42\,/\,0,50 pts \textperiodcentered\ Completo c/ ressalvas}

{\small\color{cmuted}\textit{🔵 Fez tudo, mas algo ficou incorreto ou impreciso. Esforço total, qualidade com falha.}}

\begin{itemize}[leftmargin=14pt,itemsep=1pt,topsep=3pt,parsep=0pt]
  \item[\textcolor{cgood}{$\checkmark$}] {\small\textcolor{cgood}{Dados secundários relevantes e atuais sobre o ODS escolhido}}
  \item[\textcolor{cgood}{$\checkmark$}] {\small\textcolor{cgood}{Fontes confiáveis: ONU, IBGE, artigos acadêmicos ou jornalísticos}}
  \item[\textcolor{cgood}{$\checkmark$}] {\small\textcolor{cgood}{Contextualiza o problema com base factual}}
  \item[\textcolor{cgood}{$\checkmark$}] {\small\textcolor{cgood}{Conexão clara entre os dados e o problema}}
\end{itemize}

\subsection{Matriz de Alinhamento}

\noindent\statuswarn{0,30\,/\,0,50 pts \textperiodcentered\ Faltou pouco c/ erros}

{\small\color{cmuted}\textit{🟡 Fez bastante, mas errou partes relevantes. Metade do caminho com qualidade irregular.}}

\begin{itemize}[leftmargin=14pt,itemsep=1pt,topsep=3pt,parsep=0pt]
  \item[\textcolor{cgood}{$\checkmark$}] {\small\textcolor{cgood}{Preenchida com perguntas reais do grupo sobre o problema}}
  \item[\textcolor{cgood}{$\checkmark$}] {\small\textcolor{cgood}{Reflete o que o grupo sabe e precisa descobrir}}
  \item[\textcolor{cgood}{$\checkmark$}] {\small\textcolor{cgood}{Serviu de base para planejar a pesquisa primária}}
\end{itemize}

A matriz de alinhamento não foi preenchida corretamente ! A linguagem de algumas partes está incorreta, está no tempo verbal errado. \\
O tempo verbal realmente está um pouco misturado:

\begin{itemize}[leftmargin=14pt,itemsep=1pt,topsep=3pt]
  \item "\textbackslash\{\}textbf\{Algumas pessoas já têm\}" → presente do indicativo ✅
  \item "\textbackslash\{\}textbf\{outras que ainda não doaram\}" → pretérito perfeito ❌
  \item "\textbackslash\{\}textbf\{também têm interesse\}" → presente do indicativo ✅
  \item "\textbackslash\{\}textbf\{falta direcionamento\}" → presente do indicativo ✅
\end{itemize}
Eu não sou tão crítica a esse ponto, porém atrapalhou a minha interpretação... então eu tive que pontuar de acordo com isso...

\textbf{“O que o grupo quer descobrir?”} esteja \textbf{no presente}, refletindo a situação atual, e seja \textbf{direta ou narrativa}, sem misturar tempos verbais ou repetir ideias.\\
\textbf{Versão em tópicos (direta):}

\begin{itemize}[leftmargin=14pt,itemsep=1pt,topsep=3pt]
  \item O costume das pessoas em doar alimentos.
  \item O interesse de pessoas que ainda não doaram em realizar doações.
  \item Se existe direcionamento sobre como e onde realizar doações.
\end{itemize}
\textbf{Versão narrativa (fluida):}

\begin{itemize}[leftmargin=14pt,itemsep=1pt,topsep=3pt]
  \item "Queremos descobrir se as pessoas têm o costume de doar alimentos, se aquelas que ainda não doaram possuem interesse em doar, e se existe direcionamento sobre como e onde realizar essas doações."
\end{itemize}
Existe um tópico que não foi preenchido "o que está FORA do escopo", rtalve zpor não compreenderem, mas aqui caberia informar aquilo que vocês não tem o poder para mudar, que envovle intererencia extenr aod govenro, estado, enfim algo que vcs, como grupo, não tem como resolver. Um exemplo disso seria fomenta rmais iniciativas para doações de alimentos ou alguma lei... não sei, voc~es teriamm qeu ter pesqusaido para saber como o país trata essa questão. É inaceitavel ficar em branco. \\
\\
O tópico de "Resultado esperado da imersão está errado ! Como que atráves de opesquisa no google, formulário, pesquisa de campo voc~es esperam diminuir o deperfico de alimentos ? Sendo que tudo isso envolve ações passivas e não ativas que seria a solução de vocês... ? Nesse tópico, deveria ser falado que vcs esperam sanar suas dúvidas,e ntender a faixa/limtie da propagação de doações, enfim. \\
\\
A etapa de imersão, é algo votlado para conehcimento, coleta de dados , imersão e profundiade no problema e pesquisa. É algo altamente investigativo....\\

Algo que também notei é a repetição na forma como vocês expressam o problema da doação de alimentos e o derpecidico dele.\\
\\
Em relação ao público algo que é afetado pelo problema, acredito que faltou algumas pessoas... somente pessoas ewm sitaçã de rua e com menos condições finacneiras ? Tem certeza ? A familai brasielria média ou de classe média não é afetado por isso ? \\

\\

\subsection{Matriz CSD}

\noindent\statuswarn{0,35\,/\,0,50 pts \textperiodcentered\ Completo, inadequado}

{\small\color{cmuted}\textit{🟣 Preencheu todos os itens, porém boa parte do conteúdo não atende ao propósito da atividade.}}

\begin{itemize}[leftmargin=14pt,itemsep=1pt,topsep=3pt,parsep=0pt]
  \item[\textcolor{cgood}{$\checkmark$}] {\small\textcolor{cgood}{Os três quadrantes preenchidos: Certezas, Suposições, Dúvidas}}
  \item[\textcolor{cgood}{$\checkmark$}] {\small\textcolor{cgood}{Conteúdo específico do projeto — não genérico}}
  \item[\textcolor{cgood}{$\checkmark$}] {\small\textcolor{cgood}{Revela pensamento crítico}}
  \item[\textcolor{cgood}{$\checkmark$}] {\small\textcolor{cgood}{Coerência entre a CSD e o que foi pesquisado}}
\end{itemize}

\subsection{Pesquisa Primária}

\noindent\statuswarn{0,70\,/\,1,00 pts \textperiodcentered\ Completo, inadequado}

{\small\color{cmuted}\textit{🟣 Preencheu todos os itens, porém boa parte do conteúdo não atende ao propósito da atividade.}}

\subsection{Relatório de Imersão}

\noindent\statuswarn{0,70\,/\,1,00 pts \textperiodcentered\ Completo, inadequado}

{\small\color{cmuted}\textit{🟣 Preencheu todos os itens, porém boa parte do conteúdo não atende ao propósito da atividade.}}

\begin{itemize}[leftmargin=14pt,itemsep=1pt,topsep=3pt,parsep=0pt]
  \item[\textcolor{cgood}{$\checkmark$}] {\small\textcolor{cgood}{Estrutura formal: introdução, resumo, metodologia, destaques, referências}}
  \item[$\square$] {\small\textcolor{cmuted}{Sintetiza achados da pesquisa primária e secundária}}
  \item[\textcolor{cgood}{$\checkmark$}] {\small\textcolor{cgood}{Há interpretação e conexão com o problema}}
  \item[\textcolor{cgood}{$\checkmark$}] {\small\textcolor{cgood}{Os achados alimentam a construção da persona}}
\end{itemize}

\section{Critérios e Níveis de Avaliação}

Esta seção apresenta os níveis utilizados e as penalizações por atraso, para transparência na avaliação.

\subsection{Níveis de Completude}

\begin{itemize}[leftmargin=14pt,itemsep=2pt,topsep=3pt]
  \item \textbf{Completo} (100\%): Fez tudo e fez certo. Critério plenamente atendido.
  \item \textbf{Completo c/ ressalvas} (85\%): Fez tudo, mas algo ficou incorreto ou impreciso. Esforço total, qualidade com falha.
  \item \textbf{Completo, mas inadequado} (70\%): Preencheu todos os itens, porém boa parte do conteúdo não atende ao propósito da atividade.
  \item \textbf{Faltou pouco} (65\%): Fez bastante e acertou o que fez. Faltaram partes, mas o que entregou presta.
  \item \textbf{Faltou pouco c/ erros} (60\%): Fez bastante, mas errou partes relevantes. Metade do caminho com qualidade irregular.
  \item \textbf{Faltou muito} (40\%): Fez pouco, mas o que fez está correto. Entrega incompleta, execução ok.
  \item \textbf{Faltou muito e errou} (25\%): Fez pouco e ainda errou. Esforço baixo, qualidade baixa.
  \item \textbf{Errado} (10\%): Tentou, mas a abordagem foi incorreta por completo. Reconhece a tentativa.
  \item \textbf{Não fez} (0\%): Critério não entregue.
\end{itemize}

\subsection{Penalizações por Atraso}

\begin{itemize}[leftmargin=14pt,itemsep=2pt,topsep=3pt]
  \item \textbf{Até 1 dia}: \(-10\%\) da nota máx.
  \item \textbf{2–3 dias}: \(-20\%\) da nota máx.
  \item \textbf{Mais de 3 dias}: \(-30\%\) da nota máx.
  \item \textbf{Não entregou}: zera
\end{itemize}

\section{Devolutiva Individual}

A nota individual é calculada aplicando o fator de contribuição registrado pela professora sobre a nota do grupo em cada fase.

\subsection{Fatores de Contribuição}

\begin{itemize}[leftmargin=14pt,itemsep=1pt,topsep=3pt]
  \item \textbf{Liderou} (100\%): Coordenou e entregou
  \item \textbf{Participou} (100\%): Contribuiu ativamente
  \item \textbf{Participou pouco} (70\%): Envolvimento limitado
  \item \textbf{Só fez sua parte} (40\%): Mínimo, sem engajamento
  \item \textbf{Não participou} (0\%): Sem contribuição registrada
\end{itemize}

\subsection{Notas por Integrante}

\subsubsection*{Emily Maria}

\begin{itemize}[leftmargin=14pt,itemsep=2pt,topsep=3pt,parsep=0pt]
  \item \textbf{Design Thinking — Fase 1 — Imersão}: 2,47\,pts — \textit{Participou (100\%)}

  \item[\textbf{=}] \textbf{Total Individual: 2,47\,pts}
\end{itemize}

\subsubsection*{Mariana}

\begin{itemize}[leftmargin=14pt,itemsep=2pt,topsep=3pt,parsep=0pt]
  \item \textbf{Design Thinking — Fase 1 — Imersão}: 2,47\,pts — \textit{Participou (100\%)}

  \item[\textbf{=}] \textbf{Total Individual: 2,47\,pts}
\end{itemize}

\subsubsection*{Nedson}

\begin{itemize}[leftmargin=14pt,itemsep=2pt,topsep=3pt,parsep=0pt]
  \item \textbf{Design Thinking — Fase 1 — Imersão}: 2,47\,pts — \textit{Participou (100\%)}

  \item[\textbf{=}] \textbf{Total Individual: 2,47\,pts}
\end{itemize}


\section{Considera\c{c}\~oes Finais}

Os ajustes indicados s\~ao de natureza metodol\'{o}gica e devem ser
incorporados antes do avan\c{c}o para a pr\'{o}xima fase.
Continuem assim.

% ── Assinatura ───────────────────────────────────────────────

\noindent\rule{\linewidth}{0.4pt}

\begin{flushright}
\textbf{Profa.\ Samara Silva}\\
{\small\color{cmuted} M\'{o}dulo 1 \textperiodcentered\ ETE C\'{i}cero Dias}\\
{\small\color{cmuted} Recife, 1 de junho de 2026}
\end{flushright}

\end{document}
